\begin{mainquestion}{parallelrlc}{Current-Driven Parallel RLC Circuit}

\questionpoints Consider the current-driven parallel RLC circuit shown below. The output voltage $V_{\text{out}}$ is measured across the parallel network.

\begin{center}
\begin{circuitikz}[american]
    \draw
    (0,0) node[ground] {}
    to[I, l=$I_s$] (0,3)
    -- (5,3)
    (0,0) -- (5,0)
    (1.5,3) to[R, l=$R$] (1.5,0)
    (3,3) to[L, l=$L$] (3,0)
    (4.5,3) to[C, l=$C$] (4.5,0)
    (5,3) node[right] {$+$}
    (5,0) node[right] {$-$};
\end{circuitikz}
\end{center}

% =========================================================
\begin{partquestion}{deriveimpedance}
{Derive the impedance $Z(j\omega)$ facing the current source.  
Express your answer explicitly as the sum of a real part and an imaginary part.}

\begin{answer}[3.5in]
The total admittance $Y(j\omega)$ is the sum of the admittances of the parallel components:
\[
Y(j\omega) = \frac{1}{R} + \frac{1}{j\omega L} + j\omega C.
\]
To simplify, we can factor out $1/R$ to group the reactive terms into a dimensionless quantity:
\[
Y(j\omega) = \frac{1}{R} \left[ 1 + jR\left(\omega C - \frac{1}{\omega L}\right) \right]
= \frac{1}{R} \left[ 1 + j\left(\omega RC - \frac{R}{\omega L}\right) \right].
\]
The impedance is the reciprocal of the admittance, $Z(j\omega) = 1/Y(j\omega)$:
\[
Z(j\omega) = \frac{R}{1 + j\left(\omega RC - \frac{R}{\omega L}\right)}.
\]
To separate the real and imaginary parts, we multiply the numerator and denominator by the complex conjugate of the denominator, $1 - j(\omega RC - \frac{R}{\omega L})$:
\[
Z(j\omega) = \frac{R \left[ 1 - j\left(\omega RC - \frac{R}{\omega L}\right) \right]}{1 + \left(\omega RC - \frac{R}{\omega L}\right)^2}.
\]
Thus, the real and imaginary parts are:
\[
\Re\{Z(j\omega)\} = \frac{R}{1 + \left(\omega RC - \frac{R}{\omega L}\right)^2},
\quad
\Im\{Z(j\omega)\} = -\frac{R \left(\omega RC - \frac{R}{\omega L}\right)}{1 + \left(\omega RC - \frac{R}{\omega L}\right)^2}.
\]
\end{answer}

\pagebreak
\end{partquestion}

% =========================================================
\begin{partquestion}{maxmagnitude}
{At what frequency is the magnitude of the impedance $|Z(j\omega)|$ maximized?  
What is the phase angle at that frequency?}

\begin{answer}[3in]
Based on the derivation above, the magnitude of the impedance is:
\[
|Z(j\omega)| = \frac{R}{\sqrt{1 + \left(\omega RC - \frac{R}{\omega L}\right)^2}}.
\]
The magnitude is maximized when the denominator is minimized. This occurs when the squared term is zero:
\[
\omega RC - \frac{R}{\omega L} = 0 \implies \omega^2 = \frac{1}{LC}.
\]
Solving for $\omega$ gives the resonant frequency:
\[
\omega_0 = \frac{1}{\sqrt{LC}}.
\]
At this frequency, the imaginary part of the impedance is zero, so $Z(j\omega_0) = R$. Therefore, the phase angle is:
\[
\angle Z(j\omega_0) = 0.
\]
\end{answer}

\end{partquestion}

% =========================================================
\begin{partquestion}{extremephase}
{What is the phase angle of the impedance at zero frequency (DC) and at extremely high frequency?  
Provide an intuitive explanation based on the behavior of the inductor and capacitor in these limits.}

\begin{answer}[3.5in]
At DC ($\omega \to 0$), the inductor behaves like a short circuit while the capacitor behaves like an open circuit. The impedance is therefore dominated by the inductor:
\[
Z(j\omega) \approx j\omega L.
\]
Thus, the phase angle approaches:
\[
\angle Z(j\omega) \to +\frac{\pi}{2}.
\]

At very high frequency ($\omega \to \infty$), the capacitor behaves like a short circuit while the inductor behaves like an open circuit. The impedance is dominated by the capacitor:
\[
Z(j\omega) \approx \frac{1}{j\omega C}.
\]
This has a phase angle of:
\[
\angle Z(j\omega) \to -\frac{\pi}{2}.
\]
\end{answer}

\end{partquestion}

\end{mainquestion}
